\begin{proof}
For each fixed \(q\), the map in the statement is explicitly
\[
 \kappa_q:D\otimes\Sym^4V\longrightarrow\bigwedge^3\Sym^2V,
              \qquad \epsilon\otimes f\longmapsto
                     \iota_qj(\epsilon\otimes f).
\]
Naturality of polarization and contraction makes this map equivariant
for the stabilizer of \(q\) in \(GL(V)\).  The source carries
the determinant factor \(D\); the target has its ordinary exterior
representation.  We use that equivariance in both cases.

Suppose first that \(q\) has rank \(r\in\{1,2,3\}\).
Choose \(V=A\oplus B\), where \(B=\rad q\) and \(q|_A\)
is nondegenerate.  Every pure power \(b^4\), \(b\in B\),
is killed by contraction in
\eqref{eq:polarized-schur-embedding}, since \(q(bx)=0\) for
all \(x\in V\).  Pure powers span \(\Sym^4B\), so
\(D\otimes\Sym^4B\subset\ker\kappa_q\).

The stabilizer contains \(O(A,q|_A)\times GL(B)\), the
\(B\)-scaling torus, and all shears in
Lemma~\ref{lem:shear-descent}.  On the degree-one weight space
\[
 D\otimes A\otimes\Sym^3B
\]
the product-group representation is irreducible.  Indeed its untwisted
factors are the standard representation of the full orthogonal group
and the irreducible symmetric cube for \(GL(B)\).  The determinant
factor restricts to the product character \(\det_A\det_B\),
which does not change irreducibility.  For \(r=1\) the standard
factor is one-dimensional; for \(r=2\) its two isotropic weight
lines for \(SO_2\) are exchanged by a reflection in \(O_2\),
so neither is invariant under the full group.  The usual standard
representation is irreducible also for \(r=3\).

The restriction of \(\kappa_q\) to this space is not zero.
Choose an orthogonal basis with \(a=x_1\in A\),
\(q(a,a)=1\), and \(b=x_4\in B\), and set
\(\epsilon=x_1\wedge\cdots\wedge x_4\).
In the coefficient of \(st^3\) in
\(\iota_qj(\epsilon\otimes(sa+tb)^4)\), only the contraction
of the first factor contributes.  Thus
\[
 4\kappa_q(\epsilon\otimes ab^3)
       =(bx_2)\wedge(bx_3)\wedge b^2\ne0.
\]
The three factors are independent because multiplication by \(b\)
is injective on \(V\).  Irreducibility now gives injectivity on
the degree-one space.  If the kernel contained any other positive
\(A\)-degree, Lemma~\ref{lem:shear-descent}, with \(m=4\),
would give a nonzero kernel vector in that degree-one space, a
contradiction.  This proves \eqref{eq:singular-contraction-kernel}
for every singular rank, without a genericity condition on \(q\).

Now let \(q\) be nondegenerate.  After a change of coordinates,
\(q(x_ix_j)=\delta_{ij}\) and
\(\rho=q^{-1}=x_1^2+x_2^2+x_3^2+x_4^2\).
The harmonic decomposition is
\begin{equation}\label{eq:orthogonal-quartic-decomposition}
 D\otimes\Sym^4V=
 (D\otimes\mathcal H_4)\oplus
 (D\otimes\rho\mathcal H_2)\oplus(D\otimes\C\rho^2).
\end{equation}
Here \(\mathcal H_m\) is the kernel of the Laplacian in degree
\(m\).  To check the decomposition directly, the symmetric Cauchy
formula for \(V=U\otimes U'\) in degree four has partitions
\((4),(3,1),(2,2)\), giving \(SO_4\)-types
\([4,4],[2,2],[0,0]\) of dimensions \(25,9,1\).
Multiplication by the invariant \(\rho\) identifies the degree-two
and scalar pieces; differentiation verifies the harmonic description.
The three summands are irreducible and pairwise nonisomorphic for
\(SO_4\).  The determinant factor restricts trivially to this group.

The last summand is killed: as an \(O(V,q)\)-representation,
\(D\otimes\C\rho^2\simeq D\), and
Lemma~\ref{lem:orthogonal-exterior-cube} proves that the target of
\(\kappa_q\) has no such character.  It remains to show that the
other two restrictions are nonzero.  Put
\[
 h_4=x_1^4-6x_1^2x_2^2+x_2^4\in\mathcal H_4,
            \qquad h_2=x_1^2-x_2^2\in\mathcal H_2.
\]
In the basis triple
\(\tau=(x_1x_2)\wedge(x_1x_3)\wedge(x_1x_4)\),
formula~\eqref{eq:schur-polarization-formula} gives
\[
 [\tau]\kappa_q(\epsilon\otimes h_4)=2,
     \qquad [\tau]\kappa_q(\epsilon\otimes\rho h_2)=\frac23.
\]
For the first equality the terms \(x_1^4\) and
\(-6x_1^2x_2^2\) contribute \(1,1\).  For the second the
terms \(x_1^4,x_1^2x_3^2,x_1^2x_4^2\) contribute
\(1,-1/6,-1/6\); all other terms contribute zero.
These are coefficient identities, not numerical rank assumptions.
Each nonzero restriction from an irreducible summand is injective,
and the images of the two nonisomorphic summands cannot intersect.
The kernel is therefore precisely \(D\otimes\C\rho^2\).
Equivariance transports this conclusion back to the original form
\(q\), proving \eqref{eq:nonsingular-contraction-kernel}.
\end{proof}
